\section{The extended term}
\label{sec:6.7}
Everything so far rests on the Charter as courts read it, plus the corroborated license. The extended term adds three more departures. Appendix~\ref{sec:F} gives the argument for each and works through the cases.

\emph{Standing to defend.} Article 51 speaks only of an armed attack against a Member of the United Nations. The floor extends the emergency license to any authority that administers territory under the rule in \S\ref{sec:6.2}, Member or not, and to the forces of a people whose territory a disinterested body has found occupied.\note{
Whether a population may be defended shouldn't turn on a seat other states vote on.
}

\emph{The genocide exception.} The term lifts for force to stop a genocide under way, even where the Council is blocked, on a merits judgment of the International Court of Justice under the Genocide Convention, or on two findings of different kinds, at least one from a body that took evidence.\note{
For this exception alone, a provisional-measures order finding the claim plausible and the risk real counts as one of the two; a Council or Assembly resolution corroborates, and two votes never make the pair. The exception lifts only for the force and place the finding covers: against the forces committing the genocide, where they are committing it, the least that stops it. Force elsewhere against the state responsible, on the theory that the cost will make it stop, is outside it. A finding under this exception also runs against the belligerent it names: force found to be genocide is not force repelling an attack, so that belligerent's emergency license is withdrawn for the force and place the finding covers (Appendix~\ref{sec:F.4}).
}

\emph{The occupation exception.} A finding that territory is occupied gives its population's forces the defensive license above, and no more. A finding that the occupation is unlawful and must end licenses force to end it: against the occupier's forces, inside the occupied territory, the least that ends it there, and nothing that changes who lives there.\note{
The evidence is the genocide exception's: an opinion or judgment of the Court, or two findings of different kinds, one of them from a body that took evidence. Cyprus, whose north Turkey has held since 1974, and Nagorno-Karabakh, an Armenian-majority enclave inside Azerbaijan that Armenian forces held from 1994 to 2023, are the cases (Appendix~\ref{sec:F.5}).
}

Standing to defend and the two exceptions are departures from the instruments, argued to the bar every departure must clear.

\emph{What this doesn't do.} The two-kinds rule fails East Timor, which Indonesia invaded in 1975 and held until 1999, and the Golan Heights, which Israel took from Syria in 1967 and annexed in 1981.\note{
Both had the Council's word and the Assembly's but no court or commission that took evidence, so neither occupation licenses anything; two votes making a pair would license war on bloc majorities. Karabakh shows the cost when findings don't arrive (Appendix~\ref{sec:F.5}): a license that should have lapsed in 2023 didn't, and the person-directed terms were all that saw the people the exception couldn't. And a term lifts on the facts it is given, so no drafting catches a manufactured incident, and a deployment licensed on the defender's sensors alone has accepted Gleiwitz.
}
